% !TEX root = main_without_quantization.tex
\chapter*{General Introduction}
\addcontentsline{toc}{chapter}{General Introduction}
\markboth{General Introduction}{General Introduction}
\label{chap:introduction}

The rapid progression of deep learning over the past decade has positioned large-scale
neural networks as the dominant approach for a wide range of artificial intelligence tasks,
from natural language understanding to code generation and complex reasoning.
At the center of this progression stands the Transformer architecture, whose ability to model
long-range dependencies has driven sustained improvements in performance across benchmarks.
The models that result from training at this scale are, however, extremely costly to deploy.
Their memory footprint, inference latency, and energy consumption exceed the constraints
of most production environments, and make operation on resource-limited devices impractical
without significant reduction of the model before use.

Model compression has emerged as the primary response to this gap between training-time
capability and deployment-time feasibility.
Several complementary directions have been developed: reducing the numerical precision of
stored parameters, constraining the degrees of freedom used during adaptation, and removing
parameters or computational units from the network entirely.
Each direction can reduce the resource cost of running a model, and research in each area
has produced measurable gains in efficiency with limited loss in accuracy.

Despite this progress, a fundamental difficulty remains.
Most compression approaches arrive at their decisions by evaluating each removable unit
independently, scoring it according to some local criterion before the effect of the full
set of changes on the compressed model is known.
This locality is a practical convenience but not a reliable basis for high-ratio compression:
a unit that appears unimportant in isolation may carry significant weight once the units
around it have also been removed, and the way the compression budget is allocated across
the network has a large effect on the final model quality that no local score can anticipate.

The objective of this work is to address this allocation problem directly.
Rather than scoring units one at a time and aggregating the results, the proposed framework
for compression treats the configuration of the full network as the object to be optimized,
searching over complete allocation profiles under a fixed computational budget.
The goal is to find configurations that perform well as a whole, evaluated before costly
recovery, and to do so in a way that preserves diversity across structurally different
candidates and scales to the size of modern Transformer models.

This report is divided into two parts. The first part covers the state of the art
and is composed of three chapters. Chapter~1 establishes the conceptual foundations
required to follow the rest of the work, covering neural network architectures,
training and optimization, the Transformer model, hardware efficiency metrics,
and the evolutionary and multi-objective optimization methods the contribution draws on.
Chapter~2 surveys model compression, presenting the formal setup of the principal
compression families and the structural distinctions between them.
Chapter~3 reviews the published literature within each compression family,
examines the empirical evidence for their behavior, and analyzes the allocation
limitation that motivates the proposed framework for compression.

The second part presents the contribution and is composed of two chapters.
Chapter~4 describes the proposed framework for compression in full: the parameterization
of the structural units subject to search, the budget model, the surrogate evaluation metric,
the search algorithm, and the recovery procedure applied to the selected architecture.
Chapter~5 reports the experimental results and compares the framework against
relevant baselines under controlled compression budgets.

\clearpage
